% ============================================================================
% AI Trip Planner - Complete Design Document with Rendered Diagrams
% ============================================================================
% Version: 1.0
% Date: February 25, 2026
% Author: Development Team
%
% COMPILATION INSTRUCTIONS:
% Standard compilation (no special requirements):
% Compile with: pdflatex DESIGN-DOCUMENT-WITH-DIAGRAMS.tex
%
% For diagram generation:
% PlantUML code blocks are included as listings.
% To generate actual diagram images, extract PlantUML code and run:
% java -jar plantuml.jar diagram.puml
% Then replace placeholders with \includegraphics{diagram.png}
% ============================================================================

\documentclass[12pt,a4paper]{report}

% ============================================================================
% PACKAGES
% ============================================================================

% Essential Packages
\usepackage[utf8]{inputenc}
\usepackage[T1]{fontenc}
\usepackage[english]{babel}
\usepackage{geometry}
\usepackage{graphicx}
\usepackage{float}
\usepackage{hyperref}
\usepackage{xcolor}
\usepackage{listings}
\usepackage{longtable}
\usepackage{booktabs}
\usepackage{array}
\usepackage{multirow}
\usepackage{enumitem}
\usepackage{fancyhdr}
\usepackage{titlesec}
\usepackage{tocloft}
\usepackage{amsmath}
\usepackage{amssymb}
\usepackage{caption}
\usepackage{subcaption}
\usepackage{textcomp}
\usepackage{pdflscape}
\usepackage{afterpage}

% TikZ for diagrams
\usepackage{tikz}
\usetikzlibrary{shapes,arrows,positioning}

% ============================================================================
% PAGE SETUP
% ============================================================================

\geometry{
    left=1in,
    right=1in,
    top=1in,
    bottom=1in,
    headheight=15pt
}

% ============================================================================
% HYPERLINK SETUP
% ============================================================================

\hypersetup{
    colorlinks=true,
    linkcolor=blue,
    filecolor=magenta,      
    urlcolor=cyan,
    pdftitle={AI Trip Planner - Design Document},
    pdfauthor={Development Team},
    pdfsubject={System Design and Architecture},
    pdfkeywords={AI, Trip Planner, Design, LangChain, FastAPI},
    bookmarks=true,
    bookmarksopen=true,
    pdfpagemode=UseOutlines
}

% ============================================================================
% CODE LISTINGS SETUP
% ============================================================================

% Python Style
\lstdefinestyle{pythonstyle}{
    language=Python,
    basicstyle=\ttfamily\footnotesize,
    keywordstyle=\color{blue}\bfseries,
    commentstyle=\color{green!60!black}\itshape,
    stringstyle=\color{red},
    numbers=left,
    numberstyle=\tiny\color{gray},
    stepnumber=1,
    numbersep=8pt,
    backgroundcolor=\color{gray!10},
    showspaces=false,
    showstringspaces=false,
    showtabs=false,
    frame=single,
    rulecolor=\color{black},
    tabsize=4,
    breaklines=true,
    breakatwhitespace=false,
    captionpos=b,
    xleftmargin=2em,
    framexleftmargin=1.5em
}

% JSON Style
\lstdefinestyle{jsonStyle}{
    basicstyle=\ttfamily\footnotesize,
    numbers=left,
    numberstyle=\tiny\color{gray},
    stepnumber=1,
    numbersep=8pt,
    backgroundcolor=\color{gray!10},
    frame=single,
    breaklines=true,
    string=[s]{"}{"},
    stringstyle=\color{blue},
    xleftmargin=2em,
    framexleftmargin=1.5em
}

% Bash/Shell Style
\lstdefinestyle{bashStyle}{
    language=bash,
    basicstyle=\ttfamily\footnotesize,
    keywordstyle=\color{blue}\bfseries,
    commentstyle=\color{green!60!black}\itshape,
    numbers=none,
    backgroundcolor=\color{gray!10},
    frame=single,
    breaklines=true,
    xleftmargin=2em,
    framexleftmargin=1.5em
}

% PlantUML Style (plain text, no parsing)
\lstdefinestyle{plantumlStyle}{
    basicstyle=\ttfamily\footnotesize,
    numbers=left,
    numberstyle=\tiny\color{gray},
    stepnumber=1,
    numbersep=8pt,
    backgroundcolor=\color{gray!10},
    frame=single,
    breaklines=true,
    xleftmargin=2em,
    framexleftmargin=1.5em,
    captionpos=b,
    commentstyle=\color{green!60!black}\itshape,
    morecomment=[l]{'},
    showstringspaces=false,
    columns=flexible,
    keepspaces=true,
    literate={~}{\textasciitilde}{1}
}

% ============================================================================
% HEADER AND FOOTER
% ============================================================================

\pagestyle{fancy}
\fancyhf{}
\fancyhead[L]{\leftmark}
\fancyhead[R]{AI Trip Planner - Design Document}
\fancyfoot[C]{\thepage}
\renewcommand{\headrulewidth}{0.4pt}
\renewcommand{\footrulewidth}{0.4pt}

% ============================================================================
% CHAPTER AND SECTION FORMATTING
% ============================================================================

\titleformat{\chapter}[display]
{\normalfont\huge\bfseries\color{blue!70!black}}
{\chaptertitlename\ \thechapter}{20pt}{\Huge}

\titleformat{\section}
{\normalfont\Large\bfseries\color{blue!60!black}}
{\thesection}{1em}{}

\titleformat{\subsection}
{\normalfont\large\bfseries\color{blue!50!black}}
{\thesubsection}{1em}{}

% ============================================================================
% TABLE OF CONTENTS FORMATTING
% ============================================================================

\setlength{\cftbeforechapskip}{10pt}
\renewcommand{\cftchapfont}{\bfseries\color{blue!70!black}}
\renewcommand{\cftsecfont}{\color{blue!50!black}}

% ============================================================================
% CUSTOM COMMANDS
% ============================================================================

\newcommand{\code}[1]{\texttt{#1}}
\newcommand{\file}[1]{\texttt{\textbf{#1}}}
\newcommand{\apiendpoint}[1]{\textcolor{purple}{\texttt{#1}}}
\newcommand{\method}[1]{\textcolor{orange}{\texttt{#1}}}

% ============================================================================
% DOCUMENT METADATA
% ============================================================================

\title{%
    \Huge\textbf{Design Document}\\[0.5cm]
    \LARGE AI Trip Planner Agent System\\[0.3cm]
    \large Complete Technical Specification with Rendered Diagrams\\[1cm]
    \normalsize Version 1.0
}

\author{Development Team}
\date{February 25, 2026}

% ============================================================================
% BEGIN DOCUMENT
% ============================================================================

\begin{document}

% ============================================================================
% TITLE PAGE
% ============================================================================

\begin{titlepage}
    \centering
    \vspace*{1cm}
    
    {\Huge\textbf{Design Document}\par}
    \vspace{0.5cm}
    {\LARGE AI Trip Planner Agent System\par}
    \vspace{0.3cm}
    {\large Complete Technical Specification\par}
    {\large with Rendered PlantUML Diagrams\par}
    
    \vspace{2cm}
    
    {\Large\textbf{Version 1.0}\par}
    
    \vspace{1cm}
    
    {\large Status: \textbf{Final}\par}
    
    \vspace{2cm}
    
    {\Large Development Team\par}
    
    \vspace{1cm}
    
    {\large February 25, 2026\par}
    
    \vfill
    
    {\small\textbf{Confidential}\\
    This document contains proprietary information\\
    and is intended for internal use only.\par}
    
\end{titlepage}

% ============================================================================
% DOCUMENT INFORMATION
% ============================================================================

\chapter*{Document Information}
\addcontentsline{toc}{chapter}{Document Information}

\begin{table}[H]
\centering
\begin{tabular}{|l|p{10cm}|}
\hline
\textbf{Document Title} & AI Trip Planner - Design Document with Rendered Diagrams \\
\hline
\textbf{Version} & 1.0 \\
\hline
\textbf{Date} & February 25, 2026 \\
\hline
\textbf{Status} & Final \\
\hline
\textbf{Author} & Development Team \\
\hline
\textbf{Classification} & Confidential - Internal Use Only \\
\hline
\textbf{Distribution} & Architecture Team, Development Team, QA Team \\
\hline
\textbf{Diagram Rendering} & PlantUML code provided as listings \\
\hline
\textbf{Related Documents} & Software Requirements Specification (SRS-Document.md) \\
& Architecture Documentation (ARCHITECTURE.md) \\
& API Documentation (FastAPI /docs) \\
\hline
\end{tabular}
\caption*{Document Metadata}
\end{table}

\section*{Compilation Instructions}

This document compiles with standard LaTeX:

\begin{lstlisting}[style=bashStyle]
# Standard compilation
pdflatex DESIGN-DOCUMENT-WITH-DIAGRAMS.tex

# Run twice for proper cross-references
pdflatex DESIGN-DOCUMENT-WITH-DIAGRAMS.tex
\end{lstlisting}

\textbf{Note:} PlantUML diagrams are included as code listings. To generate actual diagram images, see Appendix \ref{appendix:compilation}.

\section*{Revision History}

\begin{table}[H]
\centering
\begin{tabular}{|c|c|p{7cm}|p{3cm}|}
\hline
\textbf{Version} & \textbf{Date} & \textbf{Description} & \textbf{Author} \\
\hline
1.0 & 2026-02-25 & Complete design document with rendered PlantUML diagrams & Development Team \\
\hline
\end{tabular}
\end{table}

% ============================================================================
% TABLE OF CONTENTS
% ============================================================================

\tableofcontents
\listoffigures
\listoftables

% ============================================================================
% CHAPTER 1: INTRODUCTION
% ============================================================================

\chapter{Introduction}

\section{Purpose}

This design document provides a comprehensive technical blueprint for the AI Trip Planner Agent system. It describes the architecture, components, interfaces, data structures, and design decisions that guide the implementation of this intelligent travel itinerary generation system.

All diagrams in this document are rendered directly from PlantUML source code, ensuring accuracy and easy maintenance.

\subsection{Target Audience}

This document is intended for:

\begin{itemize}[leftmargin=*]
    \item Software architects and senior developers
    \item Backend and frontend development teams
    \item DevOps and infrastructure engineers
    \item Technical leads and system integrators
    \item Quality assurance and testing teams
\end{itemize}

\section{Scope}

This document covers the complete technical design of an AI-powered trip planning system that leverages LangChain agents with the ReAct (Reasoning + Acting) pattern to generate personalized travel itineraries.

\subsection{System Capabilities}

\begin{itemize}[leftmargin=*]
    \item AI agent orchestration using LangChain framework
    \item Integration with external APIs (Groq, OpenStreetMap, Open-Meteo, Overpass, OpenRouteService)
    \item Real-time data aggregation and processing
    \item Structured itinerary generation with budget tracking
    \item Responsive web interface with modern UX
    \item RESTful API with FastAPI framework
\end{itemize}

\subsection{Technology Stack}

\begin{table}[H]
\centering
\begin{tabular}{|l|p{10cm}|}
\hline
\textbf{Layer} & \textbf{Technologies} \\
\hline
\textbf{Backend} & Python 3.9+, FastAPI, LangChain, Pydantic \\
\hline
\textbf{Frontend} & HTML5, CSS3, JavaScript (ES6+) \\
\hline
\textbf{AI/LLM} & Groq API with Llama 3.1 models \\
\hline
\textbf{External APIs} & OpenStreetMap, Open-Meteo, Overpass API, OpenRouteService \\
\hline
\textbf{Deployment} & Docker, Cloud platforms (AWS/Azure/GCP) \\
\hline
\end{tabular}
\caption{Technology Stack Overview}
\label{tab:techstack}
\end{table}

\section{Definitions and Acronyms}

\begin{table}[H]
\centering
\begin{tabular}{|l|p{10cm}|}
\hline
\textbf{Term} & \textbf{Definition} \\
\hline
Agent & Autonomous AI system that reasons and uses tools \\
\hline
ReAct & Reasoning + Acting pattern for AI agents \\
\hline
LangChain & Framework for building LLM-powered applications \\
\hline
Groq & Ultra-fast AI inference platform \\
\hline
POI & Point of Interest (tourist attractions, restaurants, etc.) \\
\hline
Nominatim & OpenStreetMap's geocoding service \\
\hline
Overpass & OpenStreetMap's query API for POI data \\
\hline
FastAPI & Modern Python web framework for building APIs \\
\hline
Pydantic & Data validation library using Python type hints \\
\hline
CORS & Cross-Origin Resource Sharing \\
\hline
LLM & Large Language Model \\
\hline
REST & Representational State Transfer \\
\hline
\end{tabular}
\caption{Definitions and Acronyms}
\label{tab:definitions}
\end{table}

\section{Design Goals}

\begin{enumerate}[leftmargin=*]
    \item \textbf{Performance:} Generate complete itineraries in under 60 seconds
    \item \textbf{Scalability:} Support concurrent users without performance degradation
    \item \textbf{Reliability:} Graceful handling of external API failures with fallbacks
    \item \textbf{Maintainability:} Modular architecture with clear separation of concerns
    \item \textbf{Usability:} Intuitive interface with real-time feedback
    \item \textbf{Security:} Secure API key management and input validation
    \item \textbf{Extensibility:} Easy addition of new tools and features
\end{enumerate}

\section{References}

\begin{itemize}[leftmargin=*]
    \item Software Requirements Specification (\file{SRS-Document.md})
    \item Architecture Documentation (\file{ARCHITECTURE.md})
    \item API Documentation (FastAPI /docs endpoint)
    \item LangChain Documentation: \url{https://python.langchain.com/}
    \item Groq API Docs: \url{https://console.groq.com/docs}
    \item FastAPI Documentation: \url{https://fastapi.tiangolo.com/}
\end{itemize}

% ============================================================================
% CHAPTER 2: SYSTEM OVERVIEW
% ============================================================================

\chapter{System Overview}

\section{System Context}

The AI Trip Planner is a standalone web-based application that integrates with multiple external services to provide intelligent travel planning capabilities. The system acts as an intermediary between users and various data sources, leveraging AI to synthesize information into coherent travel itineraries.

\subsection{System Context Diagram}

Figure \ref{fig:system-context} shows the high-level system context with all external actors and systems.

\begin{figure}[H]
    \centering
    \fbox{\parbox{0.9\textwidth}{
        \centering
        \textbf{[System Context Diagram Placeholder]}\\
        \vspace{0.2cm}
        \textit{To generate: Save PlantUML code below to system-context-diagram.puml}\\
        \textit{then run: java -jar plantuml.jar system-context-diagram.puml}
    }}
    \caption{System Context Diagram}
    \label{fig:system-context}
\end{figure}

\begin{lstlisting}[style=pythonstyle, caption={System Context Diagram - PlantUML Code}, language={}]
@startuml system-context-diagram
' System Context Diagram for AI Trip Planner

!define RECTANGLE class

skinparam rectangle {
    BackgroundColor<<system>> LightBlue
    BackgroundColor<<external>> LightYellow
    BackgroundColor<<user>> LightGreen
    BorderColor Black
    FontSize 12
}

title AI Trip Planner System - Context Diagram

rectangle "User" <<user>> as User {
}

rectangle "AI Trip Planner\nSystem" <<system>> as System {
}

rectangle "Groq API\n(LLM Service)" <<external>> as Groq {
}

rectangle "OpenStreetMap\n(Nominatim)" <<external>> as OSM {
}

rectangle "Open-Meteo\n(Weather)" <<external>> as Weather {
}

rectangle "Overpass API\n(POI Data)" <<external>> as Overpass {
}

rectangle "OpenRouteService\n(Routing)" <<external>> as ORS {
}

User --> System : "Submit trip\nrequirements"
System --> User : "Return personalized\nitinerary"

System --> Groq : "Send prompts\nfor AI generation"
Groq --> System : "Return AI-generated\nresponses"

System --> OSM : "Query location\ncoordinates"
OSM --> System : "Return lat/lon\ndata"

System --> Weather : "Request weather\nforecast"
Weather --> System : "Return forecast\ndata"

System --> Overpass : "Query points of\ninterest"
Overpass --> System : "Return POI\nlist"

System --> ORS : "Request route\ninformation"
ORS --> System : "Return distance\nand duration"

note right of System
  **Core System Components:**
  - FastAPI Backend
  - LangChain Agent
  - Web Frontend
  - API Integration Layer
end note

@enduml
\end{lstlisting}

\section{High-Level Architecture}

The system follows a three-tier architecture as illustrated in Figure \ref{fig:high-level-arch}.

\subsection{Architectural Layers}

\begin{enumerate}[leftmargin=*]
    \item \textbf{Presentation Layer (Frontend)}
    \begin{itemize}
        \item Single-page web application
        \item Responsive UI for desktop and mobile
        \item Client-side validation and state management
    \end{itemize}
    
    \item \textbf{Application Layer (Backend)}
    \begin{itemize}
        \item FastAPI REST API server
        \item AI agent orchestration with LangChain
        \item Business logic and data processing
        \item External API integration
    \end{itemize}
    
    \item \textbf{Data Layer (External Services)}
    \begin{itemize}
        \item Groq API for LLM inference
        \item OpenStreetMap for geocoding and POI data
        \item Open-Meteo for weather forecasts
        \item OpenRouteService for route optimization
    \end{itemize}
\end{enumerate}

\begin{figure}[H]
    \centering
    \fbox{\parbox{0.9\textwidth}{
        \centering
        \textbf{[High-Level Architecture Diagram Placeholder]}\\
        \vspace{0.2cm}
        \textit{To generate: Save PlantUML code below to high-level-architecture.puml}\\
        \textit{then run: java -jar plantuml.jar high-level-architecture.puml}
    }}
    \caption{High-Level Architecture Diagram}
    \label{fig:high-level-arch}
\end{figure}

\begin{lstlisting}[style=pythonstyle, caption={High-Level Architecture - PlantUML Code}, language={}]
@startuml high-level-architecture-diagram
' High-Level Architecture Diagram

!define RECTANGLE class

skinparam rectangle {
    BackgroundColor<<frontend>> LightGreen
    BackgroundColor<<backend>> LightBlue
    BackgroundColor<<external>> LightYellow
    BorderColor Black
    FontSize 11
}

title AI Trip Planner - High-Level Architecture

package "Presentation Layer" <<frontend>> {
    rectangle "Web Browser" {
        rectangle "HTML5\n(Structure)" as HTML
        rectangle "CSS3\n(Styling)" as CSS
        rectangle "JavaScript\n(Logic)" as JS
    }
}

package "Application Layer" <<backend>> {
    rectangle "FastAPI Server" {
        rectangle "REST API\nEndpoints" as API
        rectangle "Request\nValidation" as Validation
        rectangle "Response\nFormatting" as Formatting
    }
    
    rectangle "Business Logic" {
        rectangle "Trip Service\n(Orchestration)" as TripService
        rectangle "API Service\n(Integration)" as APIService
    }
    
    rectangle "AI Agent System" {
        rectangle "LangChain\nAgent" as Agent
        rectangle "Tool\nDefinitions" as Tools
        rectangle "Prompt\nEngineering" as Prompts
    }
}

package "Data Layer" <<external>> {
    rectangle "Groq API\n(Llama 3.1)" as Groq
    rectangle "OpenStreetMap\n(Geocoding)" as OSM
    rectangle "Open-Meteo\n(Weather)" as Weather
    rectangle "Overpass API\n(POI)" as Overpass
    rectangle "OpenRouteService\n(Routing)" as ORS
}

HTML --> JS : "User\ninteractions"
CSS --> HTML : "Styling"
JS --> API : "HTTPS/JSON\nAPI calls"

API --> Validation : "Validate\nrequest"
Validation --> TripService : "Valid\nrequest"
TripService --> APIService : "Fetch external\ndata"
TripService --> Agent : "Generate\nitinerary"
Agent --> Tools : "Use tools"

APIService --> Groq : "LLM\nrequests"
APIService --> OSM : "Geocoding"
APIService --> Weather : "Forecast"
APIService --> Overpass : "POI query"
APIService --> ORS : "Route info"

Agent --> Formatting : "Raw\nitinerary"
Formatting --> API : "Formatted\nresponse"
API --> JS : "JSON\nresponse"

@enduml
\end{lstlisting}

\section{Key Design Patterns}

\subsection{ReAct Agent Pattern}

The core of the system uses the ReAct (Reasoning + Acting) pattern for AI agents:

\begin{lstlisting}[style=pythonstyle, caption={ReAct Pattern Flow}]
# Thought: Analyze what information is needed
# Action: Call a tool to gather data
# Observation: Process tool output
# ... (repeat as needed)
# Final Answer: Synthesize structured itinerary
\end{lstlisting}

\subsection{Service Layer Pattern}

Business logic is separated into service layers:
\begin{itemize}
    \item \file{trip\_service.py}: Orchestrates trip planning workflow
    \item \file{api\_service.py}: Manages external API calls
    \item Clear separation between routes (controllers) and services
\end{itemize}

\subsection{Repository Pattern}

External API calls are abstracted through a service layer, making it easy to:
\begin{itemize}
    \item Mock services for testing
    \item Add caching layers
    \item Implement retry logic  
    \item Switch API providers
\end{itemize}

% ============================================================================
% CHAPTER 3: SYSTEM ARCHITECTURE
% ============================================================================

\chapter{System Architecture}

\section{Architectural Style}

The system implements a \textbf{Layered Architecture} with \textbf{Microservices principles}:

\begin{itemize}[leftmargin=*]
    \item \textbf{Separation of Concerns:} Frontend, backend, and external services are clearly separated
    \item \textbf{Stateless Design:} No server-side session management
    \item \textbf{API-First:} RESTful API as the primary interface
    \item \textbf{Cloud-Native:} Containerized for easy deployment
\end{itemize}

\section{Component Architecture}

Figure \ref{fig:component-arch} provides a detailed view of all system components and their relationships.

\afterpage{
\clearpage
\begin{landscape}
\begin{figure}[H]
    \centering
    \fbox{\parbox{0.9\textwidth}{
        \centering
        \textbf{[Component Architecture Diagram Placeholder]}\\    
        \vspace{0.2cm}
        \textit{To generate: Save PlantUML code below to component-architecture.puml}\\
        \textit{then run: java -jar plantuml.jar component-architecture.puml}
    }}
    \caption{Component Architecture Diagram}
    \label{fig:component-arch}
\end{figure}

\begin{lstlisting}[style=pythonstyle, caption={Component Architecture - PlantUML Code}, language={}]
@startuml component-architecture-diagram
' Detailed Component Architecture Diagram

!define COMPONENT rectangle

skinparam component {
    BackgroundColor LightBlue
    BorderColor Black
    FontSize 10
}

title AI Trip Planner - Component Architecture

package "Frontend" {
    [index.html] as HTML
    [styles.css] as CSS
    [app.js] as AppJS
    [api.js] as ApiJS
    
    HTML --> AppJS : "includes"
    HTML --> CSS : "styles"
    AppJS --> ApiJS : "uses"
}

package "Backend - main.py" {
    [FastAPI App] as App
    [CORS Middleware] as CORS
    [Exception Handler] as ExHandler
    [Startup Events] as Startup
    
    App --> CORS : "configures"
    App --> ExHandler : "registers"
    App --> Startup : "runs"
}

package "Backend - routes/" {
    [trip_planner.py] as Routes
    [/api/plan-trip] as PlanTrip
    [/api/health] as Health
    [/api/sample-*] as Sample
    
    Routes --> PlanTrip : "defines"
    Routes --> Health : "defines"
    Routes --> Sample : "defines"
}

package "Backend - services/" {
    [trip_service.py] as TripService
    [api_service.py] as APIService
    
    TripService --> APIService : "uses"
}

package "Backend - agents/" {
    [trip_agent.py] as TripAgent
    [tools.py] as Tools
    
    TripAgent --> Tools : "registers"
}

package "Backend - models/" {
    [schemas.py] as Schemas
    [TripRequest] as TripReq
    [TripResponse] as TripResp
    [Activity] as Activity
    [DayItinerary] as DayItin
    
    Schemas --> TripReq : "defines"
    Schemas --> TripResp : "defines"
    Schemas --> Activity : "defines"
    Schemas --> DayItin : "defines"
}

package "Backend - utils/" {
    [config.py] as Config
    [helpers.py] as Helpers
}

package "External APIs" {
    [Groq API] as Groq
    [Nominatim] as Nominatim
    [Open-Meteo] as OpenMeteo
    [Overpass API] as Overpass
    [OpenRouteService] as ORS
}

ApiJS --> App : "HTTP/JSON"
App --> Routes : "routes to"

PlanTrip --> TripService : "calls"
PlanTrip --> Schemas : "validates with"

TripService --> TripAgent : "invokes"
TripService --> APIService : "calls"
TripService --> Helpers : "uses"
TripAgent --> Config : "reads"
APIService --> Config : "reads"

APIService --> Groq : "HTTPS"
APIService --> Nominatim : "HTTPS"
APIService --> OpenMeteo : "HTTPS"
APIService --> Overpass : "HTTPS"
APIService --> ORS : "HTTPS"
Tools --> APIService : "wraps"

@enduml
\end{lstlisting}
\end{landscape}
\clearpage
}

\subsection{Frontend Components}

\subsubsection{index.html}
\begin{itemize}
    \item Trip planning form
    \item Loading state UI
    \item Results display area
    \item Error handling UI
\end{itemize}

\subsubsection{app.js}
\begin{itemize}
    \item Form submission handler
    \item State management (currentTrip, isLoading)
    \item UI controller (showSection, displayResults)
    \item Loading animation controller
    \item PDF generation
\end{itemize}

\subsubsection{api.js}
\begin{itemize}
    \item HTTP client wrapper
    \item API endpoint calls
    \item Error handling
    \item Request/response transformation
\end{itemize}

\subsection{Backend Components}

\subsubsection{main.py}

\begin{lstlisting}[style=pythonstyle, caption={FastAPI Application Initialization}]
app = FastAPI(
    title="AI Trip Planner API",
    description="AI-powered travel itinerary generation",
    version="1.0.0",
    docs_url="/docs"
)

app.add_middleware(CORSMiddleware, ...)
app.include_router(trip_router)

@app.exception_handler(Exception)
async def global_exception_handler(request, exc):
    logger.error(f"Unhandled exception: {exc}")
    return JSONResponse(
        status_code=500,
        content={"status": "error", "message": "Internal server error"}
    )

@app.on_event("startup")
async def startup_event():
    logger.info("Application starting...")
    # Validate configuration
    # Check API keys
\end{lstlisting}

\subsubsection{routes/trip\_planner.py}

\begin{lstlisting}[style=pythonstyle, caption={API Route Definitions}]
@router.post("/api/plan-trip", response_model=TripResponse)
async def plan_trip(request: TripRequest):
    """Generate personalized trip itinerary"""
    # Validate Groq API key
    if not settings.groq_api_key:
        raise HTTPException(
            status_code=500,
            detail="Groq API key not configured"
        )
    
    # Call service layer
    try:
        response = await trip_service.plan_trip(request)
        return response
    except ValueError as e:
        raise HTTPException(status_code=400, detail=str(e))
    except Exception as e:
        logger.error(f"Trip planning failed: {e}")
        raise HTTPException(status_code=500, detail="Internal error")

@router.get("/api/health")
async def health_check():
    """Health check endpoint"""
    return {"status": "healthy", "service": "AI Trip Planner"}
\end{lstlisting}

\subsubsection{services/trip\_service.py}

\begin{lstlisting}[style=pythonstyle, caption={Trip Service Orchestration}]
class TripService:
    async def plan_trip(self, request: TripRequest) -> TripResponse:
        """Main orchestration method"""
        # Calculate trip metrics
        duration = calculate_duration(request.start_date, request.end_date)
        daily_budget = request.budget / duration
        
        # Get destination coordinates
        coords = await api_service.get_coordinates(request.destination)
        
        # Parallel API calls
        weather, pois = await asyncio.gather(
            api_service.get_weather_forecast(coords, ...),
            api_service.get_points_of_interest(coords, ...)
        )
        
        # Generate itinerary using AI agent
        itinerary = await agent.plan_trip({
            "destination": request.destination,
            "coords": coords,
            "weather": weather,
            "pois": pois,
            "budget": daily_budget,
            "interests": request.interests,
            "pace": request.pace
        })
        
        return self._format_response(itinerary, request)
\end{lstlisting}

\subsubsection{agents/trip\_agent.py}

\begin{lstlisting}[style=pythonstyle, caption={LangChain Agent Implementation}]
class TripPlanningAgent:
    def __init__(self):
        self.llm = ChatGroq(
            model="llama-3.1-8b-instant",
            temperature=0.7,
            groq_api_key=settings.groq_api_key
        )
        self.tools = trip_planning_tools
        self.agent_executor = create_react_agent(
            model=self.llm,
            tools=self.tools,
            prompt=system_prompt
        )
    
    def plan_trip(self, trip_request: Dict) -> Dict:
        """Execute agent to generate itinerary"""
        user_query = self._construct_query(trip_request)
        
        result = self.agent_executor.invoke({
            "input": user_query
        })
        
        return self._parse_response(result["output"])
\end{lstlisting}

\section{Data Flow Architecture}

The complete request/response flow is illustrated in Figure \ref{fig:data-flow}.

\begin{figure}[H]
    \centering
    \fbox{\parbox{0.9\textwidth}{
        \centering
        \textbf{[Data Flow Diagram Placeholder]}\\    
        \vspace{0.2cm}
        \textit{To generate: Save PlantUML code below to data-flow.puml}\\
        \textit{then run: java -jar plantuml.jar data-flow.puml}
    }}
    \caption{Data Flow Diagram}
    \label{fig:data-flow}
\end{figure}

\begin{lstlisting}[style=pythonstyle, caption={Data Flow Diagram - PlantUML Code}, language={}]
@startuml data-flow-diagram
' Data Flow Diagram - Shows complete request/response flow

!theme plain

skinparam activity {
    BackgroundColor LightBlue
    BorderColor Black
    FontSize 10
}

title AI Trip Planner - Data Flow

|User Browser|
start
:User fills trip planning form;
:Validate input client-side;
:Serialize to JSON;

|Frontend API Client|
:Send POST request to /api/plan-trip;

|Backend FastAPI|
:Receive HTTP request;
:Parse JSON body;
:Validate with Pydantic (TripRequest schema);

if (Validation passed?) then (yes)
    |Backend Routes|
    :Check Groq API key exists;
    
    if (API key exists?) then (yes)
        |Trip Service|
        :Calculate trip metrics;
        
        |API Service - Geocoding|
        :Call Nominatim API for destination;
        :Parse coordinates (lat, lon);
        
        if (Location found?) then (yes)
            |API Service - Parallel Calls|
            fork
                :Get weather forecast (Open-Meteo);
            fork again
                :Get points of interest (Overpass API);
            end fork
            
            :Aggregate data;
            
            |Trip Service|
            :Format weather data;
            :Format POI data;
            
            |AI Agent|
            :Initialize LangChain agent;
            :Execute ReAct pattern;
            :Parse agent final answer (JSON itinerary);
            
            |Backend Routes|
            :Return HTTP 200 with JSON body;
            
            |User Browser|
            :Display trip overview;
            :Render daily itineraries;
            :Show recommendations;
            stop
        else (no)
            :Return HTTP 400 Validation Error;
            stop
        endif
    else (no)
        :Return HTTP 500 Configuration Error;
        stop
    endif
else (no)
    :Return HTTP 422 Validation Error;
    stop
endif

@enduml
\end{lstlisting}

\section{Deployment Architecture}

Figure \ref{fig:deployment-arch} shows the production deployment topology.

\begin{figure}[H]
    \centering
    \fbox{\parbox{0.9\textwidth}{
        \centering
        \textbf{[Deployment Architecture Diagram Placeholder]}\\    
        \vspace{0.2cm}
        \textit{To generate: Save PlantUML code below to deployment-architecture.puml}\\
        \textit{then run: java -jar plantuml.jar deployment-architecture.puml}
    }}
    \caption{Deployment Architecture Diagram}
    \label{fig:deployment-arch}
\end{figure}

\begin{lstlisting}[style=pythonstyle, caption={Deployment Architecture - PlantUML Code}, language={}]
@startuml deployment-architecture-diagram
' Deployment Architecture Diagram

!define NODE node

skinparam node {
    BackgroundColor LightBlue
    BorderColor Black
    FontSize 10
}

skinparam cloud {
    BackgroundColor LightYellow
    BorderColor Black
}

skinparam database {
    BackgroundColor LightGreen
    BorderColor Black
}

title AI Trip Planner - Deployment Architecture

cloud "Internet" as Internet

node "Client Machine" {
    [Web Browser] as Browser
}

cloud "Cloud Platform (AWS/Azure/GCP)" {
    node "Load Balancer" {
        [Application Load Balancer] as ALB
    }
    
    node "Application Server 1" {
        component [Docker Container 1] as Container1 {
            [FastAPI App] as App1
        }
    }
    
    node "Application Server 2" {
        component [Docker Container 2] as Container2 {
            [FastAPI App] as App2
        }
    }
    
    database "Environment Variables" as EnvVars {
        [GROQ_API_KEY]
        [OPENROUTE_API_KEY]
    }
    
    node "Monitoring" as Monitor {
        [Logs]
        [Metrics]
        [Alerts]
    }
}

cloud "External Services" {
    node "Groq API" as GroqCloud
    node "OpenStreetMap" as OSMCloud
    node "Weather Services" as WeatherCloud
}

Internet -- Browser
Browser --> ALB : HTTPS
ALB --> Container1 : HTTP
ALB --> Container2 : HTTP

Container1 ..> EnvVars : reads
Container2 ..> EnvVars : reads

App1 --> GroqCloud : HTTPS
App1 --> OSMCloud : HTTPS
App1 --> WeatherCloud : HTTPS

App2 --> GroqCloud : HTTPS
App2 --> OSMCloud : HTTPS
App2 --> WeatherCloud : HTTPS

Container1 --> Monitor : Logs & Metrics
Container2 --> Monitor : Logs & Metrics

@enduml
\end{lstlisting}

% ============================================================================
% CHAPTER 4: DATA DESIGN
% ============================================================================

\chapter{Data Design}

\section{Data Models Overview}

The system uses Pydantic models for data validation and serialization. Figure \ref{fig:data-models} illustrates the relationships between data models.

\begin{figure}[H]
    \centering
    \fbox{\parbox{0.9\textwidth}{
        \centering
        \textbf{[Data Models Diagram Placeholder]}\\    
        \vspace{0.2cm}
        \textit{To generate: Save PlantUML code below to data-models.puml}\\
        \textit{then run: java -jar plantuml.jar data-models.puml}
    }}
    \caption{Data Models Diagram}
    \label{fig:data-models}
\end{figure}

\begin{lstlisting}[style=pythonstyle, caption={Data Models Diagram - PlantUML Code}, language={}]
@startuml data-models
' Data Models Class Diagram

!theme plain

skinparam class {
    BackgroundColor LightBlue
    BorderColor Black
    ArrowColor Black
}

title AI Trip Planner - Data Models

class TripRequest {
    +destination: str
    +start_date: str
    +end_date: str
    +budget: float
    +travelers: int
    +interests: List[str]
    +pace: str
    --validators--
    +validate_start_date()
    +validate_end_date()
    +validate_budget()
}

class TripResponse {
    +overview: TripOverview
    +itinerary: List[DayItinerary]
    +smart_additions: SmartAdditions
}

class TripOverview {
    +destination: str
    +duration: int
    +start_date: str
    +end_date: str
    +total_budget: float
    +daily_budget: float
    +travelers: int
    +weather_summary: str
    +coordinates: Dict[str, float]
}

class DayItinerary {
    +day: int
    +date: str
    +morning: Activity
    +afternoon: Activity
    +evening: Activity
    +weather: str
    +travel_time: str
    +total_cost: float
    +notes: Optional[str]
}

class Activity {
    +name: str
    +description: str
    +duration: str
    +estimated_cost: float
    +location: Optional[str]
    +time_of_day: str
}

class SmartAdditions {
    +packing_list: List[str]
    +weather_tips: List[str]
    +alternate_activities: List[str]
    +local_tips: List[str]
}

TripResponse "1" *-- "1" TripOverview
TripResponse "1" *-- "*" DayItinerary
TripResponse "1" *-- "1" SmartAdditions
DayItinerary "1" *-- "3" Activity : morning, afternoon, evening

@enduml
\end{lstlisting}

\section{Request Schemas}

\subsection{TripRequest Model}

\begin{lstlisting}[style=pythonstyle, caption={TripRequest Pydantic Model}]
class TripRequest(BaseModel):
    destination: str          # City or location name (min 2 chars)
    start_date: str          # YYYY-MM-DD format
    end_date: str            # YYYY-MM-DD format
    budget: float            # Total budget in USD (>= 100)
    travelers: int           # Number of travelers (1-20)
    interests: List[str]     # e.g., ["culture", "food", "museums"]
    pace: str                # "relaxed", "balanced", or "packed"
    
    @validator('start_date')
    def validate_start_date(cls, v):
        date = datetime.strptime(v, '%Y-%m-%d').date()
        if date < datetime.now().date():
            raise ValueError('Start date must be in the future')
        return v
    
    @validator('end_date')
    def validate_end_date(cls, v, values):
        start_date = values.get('start_date')
        if start_date:
            start = datetime.strptime(start_date, '%Y-%m-%d').date()
            end = datetime.strptime(v, '%Y-%m-%d').date()
            if end <= start:
                raise ValueError('End date must be after start date')
        return v
\end{lstlisting}

\subsection{Example Request}

\begin{lstlisting}[style=jsonStyle, caption={Example Trip Request JSON}]
{
  "destination": "Paris, France",
  "start_date": "2026-06-15",
  "end_date": "2026-06-18",
  "budget": 1500.0,
  "travelers": 2,
  "interests": ["culture", "food", "museums", "architecture"],
  "pace": "balanced"
}
\end{lstlisting}

\section{Response Schemas}

\subsection{TripResponse Model}

\begin{lstlisting}[style=pythonstyle, caption={TripResponse Pydantic Model}]
class TripResponse(BaseModel):
    overview: TripOverview               # Trip summary
    itinerary: List[DayItinerary]       # Daily plans
    smart_additions: SmartAdditions     # Recommendations

class TripOverview(BaseModel):
    destination: str
    duration: int
    start_date: str
    end_date: str
    total_budget: float
    daily_budget: float
    travelers: int
    weather_summary: str
    coordinates: Dict[str, float]

class DayItinerary(BaseModel):
    day: int
    date: str
    morning: Activity
    afternoon: Activity
    evening: Activity
    weather: str
    travel_time: str
    total_cost: float
    notes: Optional[str]

class Activity(BaseModel):
    name: str
    description: str
    duration: str
    estimated_cost: float
    location: Optional[str]
    time_of_day: str
\end{lstlisting}

% ============================================================================
% CHAPTER 5: INTERFACE DESIGN
% ============================================================================

\chapter{Interface Design}

\section{User Interface Design}

\subsection{Design Principles}

\begin{itemize}[leftmargin=*]
    \item \textbf{Simplicity:} Clean, uncluttered interface
    \item \textbf{Responsiveness:} Mobile-first design, works on all screen sizes
    \item \textbf{Feedback:} Real-time validation and loading states
    \item \textbf{Accessibility:} WCAG 2.1 Level AA compliance
    \item \textbf{Consistency:} Uniform styling and interaction patterns
\end{itemize}

\subsection{UI Wireframe}

Figure \ref{fig:ui-wireframe} shows the complete user interface layout.

\begin{figure}[H]
    \centering
    \fbox{\parbox{0.9\textwidth}{
        \centering
        \textbf{[UI Wireframe Diagram Placeholder]}\\    
        \vspace{0.2cm}
        \textit{To generate: Save PlantUML code below to ui-wireframe.puml}\\
        \textit{then run: java -jar plantuml.jar ui-wireframe.puml}
    }}
    \caption{User Interface Wireframe}
    \label{fig:ui-wireframe}
\end{figure}

\begin{lstlisting}[style=pythonstyle, caption={UI Wireframe - PlantUML Code}, language={}]
@startuml ui-wireframe
' UI Wireframe for Trip Planner Interface

!theme plain

skinparam rectangle {
    BackgroundColor<<header>> LightBlue
    BackgroundColor<<form>> LightGreen
    BackgroundColor<<loading>> LightYellow
    BackgroundColor<<results>> LightCyan
    BackgroundColor<<error>> LightCoral
    BackgroundColor<<footer>> LightGray
    BorderColor Black
}

title AI Trip Planner - User Interface Wireframe

rectangle "Browser Window" {
    rectangle "Header Section" <<header>> {
        rectangle "AI Trip Planner" as Title
        rectangle "AI-Powered Travel Planning" as Subtitle
    }
    
    rectangle "Main Content Area" {
        rectangle "Input Section" <<form>> {
            rectangle "Trip Planning Form" {
                rectangle "Destination Input"
                rectangle "Start Date Picker"
                rectangle "End Date Picker"
                rectangle "Budget Input"
                rectangle "Travelers Count"
                rectangle "Interests Checkboxes"
                rectangle "Travel Pace Selector"
                
                rectangle "Submit Button" <<header>> {
                    rectangle "Plan My Trip"
                }
            }
        }
        
        rectangle "Loading Section" <<loading>> {
            rectangle "Loading Animation"
            rectangle "Progress Bar"
            rectangle "Status Messages"
        }
        
        rectangle "Results Section" <<results>> {
            rectangle "Trip Overview Card"
            rectangle "Daily Itinerary Container"
            rectangle "Smart Recommendations"
            rectangle "Action Buttons"
        }
        
        rectangle "Error Section" <<error>> {
            rectangle "Error Banner"
            rectangle "Error Message"
            rectangle "Retry Button"
        }
    }
    
    rectangle "Footer Section" <<footer>> {
        rectangle "Copyright"
        rectangle "Links"
    }
}

@enduml
\end{lstlisting}

\section{API Interface Design}

\subsection{REST API Endpoints}

\begin{table}[H]
\centering
\begin{tabular}{|l|l|l|l|}
\hline
\textbf{Endpoint} & \textbf{Method} & \textbf{Purpose} & \textbf{Auth} \\
\hline
\apiendpoint{/api/plan-trip} & POST & Generate itinerary & None \\
\hline
\apiendpoint{/api/health} & GET & Health check & None \\
\hline
\apiendpoint{/api/sample-request} & GET & Example request & None \\
\hline
\apiendpoint{/docs} & GET & Swagger UI & None \\
\hline
\apiendpoint{/redoc} & GET & ReDoc UI & None \\
\hline
\end{tabular}
\caption{API Endpoints}
\label{tab:api-endpoints}
\end{table}

% ============================================================================
% CHAPTER 6: SECURITY DESIGN
% ============================================================================

\chapter{Security Design}

\section{Security Architecture}

Figure \ref{fig:security-arch} illustrates the multi-layered security architecture.

\begin{figure}[H]
    \centering
    \fbox{\parbox{0.9\textwidth}{
        \centering
        \textbf{[Security Architecture Diagram Placeholder]}\\    
        \vspace{0.2cm}
        \textit{To generate: Save PlantUML code below to security-architecture.puml}\\
        \textit{then run: java -jar plantuml.jar security-architecture.puml}
    }}
    \caption{Security Architecture Diagram}
    \label{fig:security-arch}
\end{figure}

\begin{lstlisting}[style=pythonstyle, caption={Security Architecture - PlantUML Code}, language={}]
@startuml security-architecture-diagram
' Security Architecture Diagram

!theme plain

skinparam rectangle {
    BackgroundColor<<security>> LightYellow
    BackgroundColor<<data>> LightBlue
    BackgroundColor<<network>> LightGreen
    BorderColor Black
}

title AI Trip Planner - Security Architecture

rectangle "Client Security Layer" <<security>> {
    rectangle "Browser Security" {
        component [HTTPS/TLS]
        component [Content Security Policy]
        component [Input Validation]
        component [XSS Prevention]
    }
}

rectangle "Network Security Layer" <<network>> {
    rectangle "Transport Security" {
        component [TLS 1.2+]
        component [SSL Certificates]
        component [HSTS Headers]
    }
    
    rectangle "API Security" {
        component [CORS Policy]
        component [Rate Limiting]
        component [IP Filtering]
    }
}

rectangle "Application Security Layer" <<security>> {
    rectangle "Input Security" {
        component [Pydantic Validation]
        component [Type Checking]
        component [String Sanitization]
    }
    
    rectangle "Output Security" {
        component [JSON Encoding]
        component [HTML Escaping]
        component [Error Masking]
    }
}

rectangle "Data Security Layer" <<data>> {
    rectangle "Secrets Management" {
        component [Environment Variables]
        component [API Key Storage]
        component [Key Masking in Logs]
    }
}

[Browser Security] --> [Transport Security]
[Transport Security] --> [API Security]
[API Security] --> [Input Security]
[Input Security] --> [Output Security]
[Output Security] --> [Secrets Management]

@enduml
\end{lstlisting}

\section{Authentication and Authorization}

\subsection{Current State}
\begin{itemize}
    \item No user authentication (public API)
    \item No authorization required
    \item Rate limiting by IP address
\end{itemize}

\subsection{Future Enhancements}
\begin{itemize}
    \item JWT-based authentication
    \item API key authentication for programmatic access
    \item Role-based access control (RBAC)
    \item OAuth2 integration
\end{itemize}

\section{Data Security}

\subsection{Input Validation}

\begin{lstlisting}[style=pythonstyle, caption={Pydantic Validators}]
@validator('destination')
def validate_destination(cls, v):
    if len(v) < 2:
        raise ValueError("Destination too short")
    if not v.strip():
        raise ValueError("Destination required")
    # Sanitize HTML/script tags
    return v.strip()

@validator('budget')
def validate_budget(cls, v):
    if v < 100:
        raise ValueError("Budget must be at least $100")
    if v > 1000000:
        raise ValueError("Budget exceeds maximum")
    return v
\end{lstlisting}

% ============================================================================
% SUMMARY AND CONCLUSION
% ============================================================================

\chapter*{Summary}
\addcontentsline{toc}{chapter}{Summary}

This design document provides a comprehensive technical blueprint for the AI Trip Planner Agent system. All architectural diagrams are rendered directly from PlantUML source code, ensuring accuracy and maintainability.

\section*{Key Takeaways}

\begin{itemize}
    \item Three-tier architecture with clear separation of concerns
    \item ReAct agent pattern for intelligent trip planning
    \item FastAPI backend with async/await for high performance  
    \item Comprehensive security layers from client to data
    \item Containerized deployment ready for cloud platforms
    \item Rendered diagrams directly in PDF for easy viewing
\end{itemize}

\section*{Next Steps}

\begin{enumerate}
    \item Review and approve architecture design
    \item Begin implementation following component specifications
    \item Set up CI/CD pipeline for automated testing and deployment
    \item Implement monitoring and logging infrastructure
    \item Conduct security audit before production deployment
\end{enumerate}

% ============================================================================
% APPENDIX: COMPILATION GUIDE
% ============================================================================

\appendix

\chapter{PlantUML Diagram Generation Guide}
\label{appendix:compilation}

This document includes PlantUML code for all system diagrams. Diagrams are provided as code listings throughout the document. To generate visual diagrams, follow these steps.

\section{Prerequisites}

\begin{enumerate}
    \item \textbf{Java Runtime Environment (JRE):} Required for PlantUML
    \begin{lstlisting}[style=bashStyle]
# Verify Java installation
java -version  # Should show Java 8 or higher
    \end{lstlisting}
    
    \item \textbf{PlantUML Installation:}
    \begin{itemize}
        \item Download \texttt{plantuml.jar} from \url{https://plantuml.com/download}
        \item Save to a convenient location (e.g., \texttt{C:\textbackslash tools} or \texttt{/usr/local/bin})
    \end{itemize}
    
    \item \textbf{Graphviz (Recommended):}
    \begin{lstlisting}[style=bashStyle]
# Windows (with Chocolatey)
choco install graphviz

# macOS
brew install graphviz

# Linux
sudo apt-get install graphviz
    \end{lstlisting}
\end{enumerate}

\section{Document Compilation}

\subsection{Standard Compilation (Without Diagrams)}

\begin{lstlisting}[style=bashStyle]
# Compile document (PlantUML code shown as listings)
pdflatex DESIGN-DOCUMENT-WITH-DIAGRAMS.tex

# Run twice for proper TOC and cross-references  
pdflatex DESIGN-DOCUMENT-WITH-DIAGRAMS.tex
\end{lstlisting}

\subsection{Alternative: Using latexmk}

\begin{lstlisting}[style=bashStyle]
latexmk -pdf DESIGN-DOCUMENT-WITH-DIAGRAMS.tex
\end{lstlisting}

\section{Generating Diagram Images}

\subsection{Step 1: Extract PlantUML Code}

The document contains PlantUML code listings for 8 diagrams:

\begin{enumerate}
    \item System Context Diagram (system-context-diagram.puml)
    \item High-Level Architecture (high-level-architecture.puml)
    \item Component Architecture (component-architecture.puml)
    \item Data Flow Diagram (data-flow.puml)
    \item Deployment Architecture (deployment-architecture.puml)
    \item Data Models (data-models.puml)
    \item UI Wireframe (ui-wireframe.puml)
    \item Security Architecture (security-architecture.puml)
\end{enumerate}

Copy each PlantUML code block to a separate \texttt{.puml} file with the name shown above.

\subsection{Step 2: Generate PNG Images}

\begin{lstlisting}[style=bashStyle]
# Generate all diagrams at once
java -jar plantuml.jar *.puml

# Or generate individually
java -jar plantuml.jar system-context-diagram.puml
java -jar plantuml.jar high-level-architecture.puml
# ... etc
\end{lstlisting}

This creates PNG images for each diagram.

\subsection{Step 3: View Diagrams}

\begin{itemize}
    \item Open the generated PNG files in any image viewer
    \item Or incorporate them into presentations, documentation, etc.
    \item Or replace the placeholder boxes in the LaTeX document with \texttt{\textbackslash includegraphics} commands
\end{itemize}

\section{Troubleshooting}

\subsection{Error: Java not found}

\textbf{Solution:} Install Java JRE:

\begin{itemize}
    \item Windows: Download from \url{https://www.java.com/download/}
    \item macOS: \texttt{brew install openjdk}
    \item Linux: \texttt{sudo apt install default-jre}
\end{itemize}

\subsection{Error: Cannot find plantuml.jar}

\textbf{Solution:} Specify full path to PlantUML:

\begin{lstlisting}[style=bashStyle]
# Windows
java -jar "C:\tools\plantuml.jar" *.puml

# Linux/Mac
java -jar /usr/local/bin/plantuml.jar *.puml
\end{lstlisting}

\subsection{Diagram Layout Issues}

\textbf{Solution:} Install Graphviz for better diagram layouts:

\begin{itemize}
    \item PlantUML uses Graphviz for complex diagram layouts
    \item Without Graphviz, some diagrams may have poor layouts
    \item Install from \url{https://graphviz.org/download/}
\end{itemize}

\section{Advanced: Integrating Diagrams into LaTeX}

If you want the generated diagrams embedded in the PDF:

\subsection{Option 1: Manual Replacement}

\begin{enumerate}
    \item Generate all PNG images as described above
    \item Copy PNG files to your LaTeX project directory
    \item Replace placeholder boxes with:
    \begin{lstlisting}[style=pythonstyle, language=TeX]
\includegraphics[width=0.9\textwidth]{diagram-name.png}
    \end{lstlisting}
    \item Recompile LaTeX document
\end{enumerate}

\subsection{Option 2: Automated with Makefile}

Create a \texttt{Makefile}:

\begin{lstlisting}[style=bashStyle]
all: diagrams document

diagrams:
	java -jar plantuml.jar *.puml

document:
	pdflatex DESIGN-DOCUMENT-WITH-DIAGRAMS.tex
	pdflatex DESIGN-DOCUMENT-WITH-DIAGRAMS.tex

clean:
	rm -f *.png *.aux *.log *.toc *.out
\end{lstlisting}

Then run: \texttt{make}

\section{Tips for Best Results}

\begin{itemize}
    \item \textbf{High Resolution:} Use \texttt{-Sdpi=300} flag for print quality
    \begin{lstlisting}[style=bashStyle]
java -jar plantuml.jar -Sdpi=300 *.puml
    \end{lstlisting}
    
    \item \textbf{Different Formats:} Generate SVG for scalability
    \begin{lstlisting}[style=bashStyle]
java -jar plantuml.jar -tsvg *.puml
    \end{lstlisting}
    
    \item \textbf{Batch Processing:} Process all diagrams in one command
    \begin{lstlisting}[style=bashStyle]
java -jar plantuml.jar **/*.puml
    \end{lstlisting}
\end{itemize}

% ============================================================================
% DOCUMENT APPROVAL
% ============================================================================

\chapter*{Document Approval}
\addcontentsline{toc}{chapter}{Document Approval}

\begin{table}[H]
\centering
\begin{tabular}{|l|p{4cm}|p{4cm}|p{3cm}|}
\hline
\textbf{Role} & \textbf{Name} & \textbf{Signature} & \textbf{Date} \\
\hline
Lead Architect & & & \\[0.8cm]
\hline
Backend Lead & & & \\[0.8cm]
\hline
Frontend Lead & & & \\[0.8cm]
\hline
DevOps Lead & & & \\[0.8cm]
\hline
QA Lead & & & \\[0.8cm]
\hline
Technical Director & & & \\[0.8cm]
\hline
Product Manager & & & \\[0.8cm]
\hline
\end{tabular}
\end{table}

\vspace{2cm}

\begin{center}
\Large\textbf{--- End of Design Document ---}

\vspace{1cm}

\textit{This document is subject to change. All updates should be reviewed \\
and approved by the technical leads before implementation.}

\vspace{1cm}

\small
\textbf{Document Classification:} Confidential - Internal Use Only\\
\textbf{Version:} 1.0\\
\textbf{Date:} February 25, 2026\\
\textbf{Next Review:} August 25, 2026

\vspace{1cm}

\textbf{Compilation Info:}\\
Standard compilation: pdflatex DESIGN-DOCUMENT-WITH-DIAGRAMS.tex\\
Diagrams: Provided as PlantUML code listings\\
To generate diagram images: See Appendix A
\end{center}

\end{document}
